\section{Managerial and Operational Interpretation}

The managerial object in this paper is the service menu, not only the final accepted ride. A platform operator controls which meeting-point and pickup-window bundles are visible, how those bundles are priced, and which eligibility filters act before passengers see the menu. The strengthened framework therefore asks managers to audit three layers separately: candidate eligibility, assortment construction, and pricing/choice response.

The exact-vs-greedy diagnostics are useful because they turn a common operational compromise into a measurable tradeoff. If greedy forward selection has a small surrogate-value gap and much lower build time, it is a defensible real-time approximation. If the gap becomes large in certain regimes, the operator can selectively raise the exact-enumeration threshold or reserve exact solving for requests with high expected value.

The robust-filtering comparison is equally important. Hard ETA filters can make the downstream assortment solver appear ineffective by removing feasible meeting-point bundles before they are evaluated. Calibrated and interval-overlap filters provide middle ground: they acknowledge ETA uncertainty without treating every noisy point estimate as a hard exclusion. The no-filter variant remains a diagnostic endpoint, useful for measuring how much menu exposure is being suppressed.

Finally, the uptake-regime comparison clarifies when menu optimization should be expected to matter. In very low-uptake settings, observed profit gaps mainly diagnose mechanism behavior because passengers rarely accept non-home bundles. In medium and higher-uptake settings, acceptance, non-home uptake, consumer surplus, and price-boundary rates become informative managerial outcomes. A platform should therefore evaluate menu policies in behaviorally live regimes before promoting one pricing rule or filtering rule as operationally superior. The cross-instance results are descriptive two-pair checks; the statistical evidence base for menu-optimization value comes from the six-pair RC benchmark and the uptake-regime stress tests.
